% Project report generated from the actual canteen_test implementation
\documentclass[12pt,a4paper]{report}
\usepackage[margin=1in]{geometry}
\usepackage{hyperref}
\usepackage{titlesec}
\usepackage{enumitem}
\usepackage{longtable}
\usepackage{booktabs}
\usepackage{setspace}
\onehalfspacing

\title{Canteen Management System: Project Report}
\author{Generated from Actual Codebase Implementation}
\date{\today}

\begin{document}

\maketitle
\tableofcontents
\clearpage

\chapter{Introduction}
\section{Purpose}
The purpose of this project report is to present a comprehensive description of the Canteen Management System built from the actual source code in the repository. This report is based directly on the implementation found in the workspace folders: \texttt{canteen-backend}, \texttt{canteen-frontend}, and \texttt{canteen-admin}.

\section{Scope}
The system provides a full-stack digital canteen solution with the following capabilities:
\begin{itemize}
  \item Mobile customer application for browsing menu items, combos, wallet management, and placing orders.
  \item Backend REST API services for authentication, menu management, order processing, wallet operations, QR-based order verification, and admin workflows.
  \item Admin dashboard for staff and management, implemented using React, React Query, Tailwind-style UI, and client-side authentication.
  \item MongoDB persistence, JWT-based authentication, and modular backend architecture.
\end{itemize}

\section{Objectives}
The primary objectives of the system are:
\begin{enumerate}
  \item Digitize canteen ordering and payment workflows.
  \item Provide a responsive mobile experience for students and staff.
  \item Support order tracking, QR-based verification, and wallet transactions.
  \item Enable administrative control over menu, combos, inventory, and orders.
\end{enumerate}

\chapter{Literature Survey}
\section{Background}
Traditional canteen systems rely on manual order taking, physical menus, and cash payments. Modern digital ordering platforms apply mobile user interfaces, backend APIs, and real-time status updates to reduce wait time and improve order accuracy.

\section{Relevant Technologies}
The implementation uses the following contemporary technologies:
\begin{itemize}
  \item FastAPI for backend REST services and automatic API documentation.
  \item MongoDB for document-oriented persistence.
  \item React Native with Expo for mobile application development.
  \item React + TypeScript + Vite for admin dashboard development.
  \item JWT authentication and HTTP Bearer token security.
\end{itemize}

\section{Comparative Context}
Compared to legacy canteen management software, this repository demonstrates a modern microservice-inspired structure with clear separation of concerns. The codebase emphasizes API modularity in the backend through separate route modules (e.g., \texttt{app/routes/menu.py}, \texttt{app/routes/order.py}, \texttt{app/routes/qr.py}) and separate frontend screens for each user flow.

\chapter{System Analysis}
\section{Stakeholders}
The main stakeholders are:
\begin{itemize}
  \item Canteen customers (students, staff) using the mobile app.
  \item Canteen managers and staff using the admin dashboard.
  \item System operators maintaining the backend and database.
\end{itemize}

\section{Functional Requirements}
Based on the actual implementation, the system supports the following functional requirements:
\begin{itemize}
  \item User registration and login with JWT authentication.
  \item Retrieve current user profile via \texttt{/api/auth/me}.
  \item Browse menu items and filter by category, trending, special, featured, and vegetarian attributes.
  \item View combo offers through \texttt{/api/combos} endpoints.
  \item Add items to a cart and place orders with \texttt{/api/orders}.
  \item Track order status using \texttt{/api/orders/{order_id}/status}.
  \item Generate and verify QR codes for order validation through \texttt{/api/qr}.
  \item Manage wallet details, set/change wallet PIN, and top up balance using \texttt{/api/wallet}.
  \item Admin operations such as menu item creation, updates, and inventory control via authenticated backend routes.
\end{itemize}

\section{Non-Functional Requirements}
The system also targets the following non-functional properties:
\begin{itemize}
  \item Modularity through separate backend route and service modules.
  \item Extensibility via clearly defined Pydantic schemas and FastAPI dependency injection.
  \item Maintainability using componentized React and React Native screens.
  \item Usability with tab navigation, drawer navigation, and contextual wallet/cart flows.
\end{itemize}

\section{System Context}
The application consists of three main subsystems:
\begin{itemize}
  \item \textbf{Backend API}: hosted by FastAPI in \texttt{canteen-backend/app/main.py}, exposing endpoints under \texttt{/api/*}.
  \item \textbf{Mobile Frontend}: implemented with Expo React Native in \texttt{canteen-frontend/App.js} and supporting screens under \texttt{canteen-frontend/src/screens}.
  \item \textbf{Admin Frontend}: implemented in \texttt{canteen-admin} using React Router, React Query, and local storage-based admin state.
\end{itemize}

\chapter{System Design}
\section{Architectural Overview}
The solution follows a client-server architecture with the following design principles:
\begin{itemize}
  \item Backend services are organized in route modules, service modules, schema definitions, and database connection utilities.
  \item The mobile client interacts with backend REST endpoints through an Axios wrapper in \texttt{canteen-frontend/src/services/api.js} and \texttt{apiService.js}.
  \item Administrative UI uses a dashboard layout with a sidebar and top navigation, controlled by an auth context in \texttt{canteen-admin/src/context/AuthContext.tsx}.
\end{itemize}

\section{Backend Design}
The backend implementation includes these key components:
\begin{itemize}
  \item \textbf{FastAPI App Initialization}: \texttt{canteen-backend/app/main.py} configures CORS, startup/shutdown events, and router inclusion.
  \item \textbf{Authentication}: \texttt{canteen-backend/app/routes/auth.py} defines registration and login endpoints, while \texttt{canteen-backend/app/core/dependencies.py} verifies JWT tokens and provides role-based dependencies.
  \item \textbf{Menu and Combo Services}: \texttt{app/routes/menu.py} and \texttt{app/routes/combo.py} expose item listing, filtering, creation, and update endpoints.
  \item \textbf{Order Management}: \texttt{app/routes/order.py} handles order placement, retrieval, status checks, and status updates.
  \item \textbf{QR Verification}: \texttt{app/routes/qr.py} supports generating order QR codes and verifying them to complete orders.
  \item \textbf{Wallet Management}: \texttt{app/routes/wallet.py} supports wallet details, transaction history, PIN management, and top-up.
\end{itemize}

\section{Frontend Design}
The mobile application architecture includes:
\begin{itemize}
  \item \textbf{Navigation}: A stack navigator in \texttt{App.js} with bottom tabs for Home, Menu, Cart, Orders, and Wallet, plus drawer navigation for profile and edit screens.
  \item \textbf{State Management}: A context provider in \texttt{canteen-frontend/src/context/CartContext.js} holds cart items, wallet balance, and user profile state.
  \item \textbf{API Layer}: \texttt{canteen-frontend/src/services/apiService.js} performs network requests with centralized error handling and response normalization.
  \item \textbf{Screen Components}: Screens such as \texttt{HomeScreen}, \texttt{MenuScreen}, \texttt{CartScreen}, \texttt{OrdersScreen}, and \texttt{WalletScreen} deliver specific user workflows.
\end{itemize}

\section{Admin Dashboard Design}
The admin dashboard uses the following design:
\begin{itemize}
  \item \textbf{Routing}: React Router and nested routes under \texttt{DashboardLayout.tsx}.
  \item \textbf{Authentication}: Admin login state is managed in \texttt{AuthContext.tsx}, persisted to local storage.
  \item \textbf{UI}: A sidebar and top navigation provide access to inventory, orders, users, and reports.
  \item \textbf{Data Fetching}: React Query is available via dependencies listed in \texttt{canteen-admin/package.json} and is intended for backend data synchronization.
\end{itemize}

\chapter{Implementation}
\section{Backend Implementation}
The repository demonstrates an implemented backend with the following actual elements:
\begin{itemize}
  \item \textbf{Startup and Shutdown}: \texttt{connect_to_mongo} and \texttt{close_mongo_connection} in \texttt{app/database/mongodb.py} are wired to FastAPI lifecycle events.
  \item \textbf{HTTP Endpoints}: All endpoints are prefixed under \texttt{/api}, such as \texttt{/api/auth}, \texttt{/api/menu}, \texttt{/api/orders}, \texttt{/api/wallet}, and \texttt{/api/qr}.
  \item \textbf{Data Validation}: Pydantic models in \texttt{canteen-backend/app/schemas} define DTOs like \texttt{RegisterRequest}, \texttt{LoginRequest}, \texttt{OrderCreate}, \texttt{MenuItemCreate}, and \texttt{WalletTopUpRequest}.
  \item \textbf{Security}: JWT decoding takes place in \texttt{app/core/dependencies.py}, and demo-hardcoded mock JWT support is present for admin dashboard integration.
\end{itemize}

\section{Mobile Frontend Implementation}
Key implementation specifics in the mobile app include:
\begin{itemize}
  \item \textbf{Root App Component}: \texttt{canteen-frontend/App.js} configures navigation and provides shared state using \texttt{CartContext}.
  \item \textbf{Screen Flow}: The app presents login/register screens, followed by a tabbed home experience containing menu, cart, orders, and wallet sections.
  \item \textbf{Device Integration}: The app uses Expo packages such as \texttt{expo-navigation-bar}, \texttt{expo-image}, and \texttt{expo-linear-gradient} for enhanced mobile UI and system navigation styling.
  \item \textbf{API Integration}: The app calls backend endpoints through Axios wrappers in \texttt{src/services/api.js} and \texttt{src/services/apiService.js}.
\end{itemize}

\section{Admin Dashboard Implementation}
Admin implementation details include:
\begin{itemize}
  \item \textbf{Auth Persistence}: \texttt{AuthProvider} stores admin token and user data in browser local storage.
  \item \textbf{Layout}: \texttt{DashboardLayout.tsx} provides a responsive wrapper containing \texttt{Sidebar} and \texttt{TopNavbar} components.
  \item \textbf{Developer Workflow}: Scripts in \texttt{canteen-admin/package.json} include \texttt{dev}, \texttt{build}, \texttt{lint}, and \texttt{preview}.
\end{itemize}

\chapter{Testing}
\section{Observed Testing Artifacts}
In the inspected codebase, there is no dedicated automated test suite or test directory present. The repository does include developer tooling that supports validation:
\begin{itemize}
  \item The frontend and admin projects expose lint scripts via \texttt{npm run lint}.
  \item The backend exposes a health endpoint at \texttt{/health} for runtime checks.
\end{itemize}

\section{Manual Verification Paths}
The actual implementation allows the following verification paths:
\begin{itemize}
  \item Run the backend with \texttt{uvicorn app.main:app --reload} and confirm the health route returns \texttt{{"status":"ok"}}.
  \item Access Swagger UI for backend API discovery and contract verification.
  \item Start Expo with \texttt{npm start} in \texttt{canteen-frontend} and validate mobile screen navigation.
  \item Start Vite with \texttt{npm run dev} in \texttt{canteen-admin} and verify dashboard pages load.
\end{itemize}

\section{Testing Limitations}
The current code snapshot does not include explicit JUnit, PyTest, or Jest test files. As a result, rigorous regression coverage is not directly observable in the repository contents.

\chapter{Results}
\section{Implemented Capabilities}
The repository demonstrates a working implementation of the following capabilities:
\begin{itemize}
  \item End-to-end user journey from registration through order placement.
  \item Backend route coverage for authentication, profile, menu, combo, order, wallet, QR verification, and inventory management.
  \item Mobile navigation structure supporting home, menu, cart, orders, wallet, profile, and QR interaction.
  \item Admin dashboard with authentication state, layout, and a foundation for inventory and orders management.
\end{itemize}

\section{Value Delivered}
The implemented system delivers:
\begin{itemize}
  \item A practical mobile-friendly ordering experience for canteen customers.
  \item A backend API surface suitable for mobile and web clients.
  \item A modular codebase that separates concerns across backend routing, business services, schema validation, and frontend UI routes.
\end{itemize}

\chapter{Conclusion}
This project report documents the actual implementation of the Canteen Management System based on the repository contents. The codebase exhibits a modern full-stack architecture, with a FastAPI backend, MongoDB persistence, an Expo React Native mobile client, and a React admin dashboard.

Future enhancements could include explicit automated tests, restored role-based access control enforcement, richer admin analytics, and production-grade deployment configuration for both backend and mobile/web clients.

\end{document}
